\section{Discussion}
\label{sec:discussion}

The evaluation indicates that TPS-ZKPI removes recoverable policy-dependent
signals from normalized public-chain artifacts. Across all PCH, PI, and SPP
tasks, supervised classifiers stay below the multi-seed permutation-test
leakage criterion, and clustering ARI is statistically indistinguishable from
zero. The gas-padding ablation (Section~\ref{sec:eval-ablation})
sharpens this conclusion: the same attacks become measurably effective once
normalization is disabled, which attributes the protection specifically to
transcript normalization rather than to weak attacks or insufficient sample
size. Together with fixed-length policy representation, reusable interpretation,
and commitment binding, these mechanisms limit the extent to which public-chain
observers infer hidden policy content or structural properties from
authorization transactions.

These results support a scoped privacy claim. The evaluation focuses on a
metadata-only public-chain adversary that observes commitments, proofs,
public parameters, calldata, events, and gas-related artifacts, but not TPS-side
witness construction or off-chain policy material. Within this scope,
transcript normalization eliminates the recoverable policy-dependent signals
we are able to measure across PCH, PI, and SPP, at a gas-cost overhead of
roughly $38$--$40\%$. We do not claim distinguishability is information-theoretically
zero; rather, no modeled attacker recovers it above the chance baseline.
Stronger adversaries, such as compromised proof services, adaptive
requester-side probing, or off-chain side channels, remain outside the current
evaluation scope and motivate the reduced-trust extensions discussed below.

The TPS trust assumption is the main architectural tradeoff. By moving
confidentiality-sensitive proof generation to the TPS, TPS-ZKPI avoids exposing
policy-bound proving material to requesters or public smart contracts, while
preserving on-chain verifiability through zero-knowledge proof verification.
However, this shifts part of the confidentiality burden to the service layer.
Future work should investigate reduced-trust TPS designs, including
distributed proof services, hardware-assisted isolation, accountable logging,
or cryptographic protocols that further reduce what any single proof service
can observe.
